\chapter{Constitution du corpus et ensembles analytiques}\label{ch:corpus}

\section{Critères de sélection}

Le corpus a été extrait de l’export de l’EDCS diffusé par le projet SDAM de l’université d’Aarhus (fichier « province Roma » du 9 décembre 2020, 93\,058 fiches pour Rome). Les critères d’inclusion, fixés par le catalogue de travail initial, sont les suivants :
\begin{enumerate}[label=(\alph*),leftmargin=2em]
\item épitaphes chrétiennes de Rome (classement « inscriptiones christianae » et « tituli sepulcrales » de la base) ;
\item datées par la base d’une année comprise entre 340 et 403 (443 inscriptions) ou d’une fourchette de cinq ans au plus (8) ;
\item auxquelles s’ajoutent 27 épitaphes dont le consulat, lisible dans le texte, n’était pas reporté dans le champ de datation de la base.
\end{enumerate}

La sélection visait à exclure les inscriptions honorifiques ou de construction et les textes datés seulement par des critères paléographiques, formulaires ou topographiques. Le contrôle montre qu’elle n’y est pas entièrement parvenue : une inscription architecturale (\nnum{321}), des actes d’achat sans commémoration et quelques textes dont l’année ne repose sur aucun consulat lisible figurent parmi les \st{nA} notices. Leur traitement est précisé ci-dessous. Le terme de 403 a été fixé sur la seule épitaphe du corpus qui mentionne la Trinité (\nnum{474}). Ce choix de borne, lié à la problématique initiale, ne correspond à aucune césure de la documentation ; il doit être lu comme une limite de travail.

\section{Notices, textes, unités documentaires}

Trois unités de compte sont distinguées.

\begin{description}[style=nextline,leftmargin=1.5em,font=\normalfont\scshape]
\item[Ensemble A — notices (\st{nA})] Ensemble documentaire initial. Les numéros 1–\st{nA} sont des identifiants stables : aucune notice n’est supprimée ni renumérotée. Une notice fusionnée reste au catalogue, avec un renvoi.
\item[Ensemble B — textes distincts datés de 340 à 403 (\st{nB})] Ensemble A moins les trois notices reconnues comme doublons d’une autre notice (106, 282, 329), le second fragment d’une plaque dont le premier fragment est déjà compté (394, joint au 393), et la notice 98, datée de 461 ou de 482. C’est l’ensemble des textes du corpus.
\item[Ensemble C — textes funéraires (\st{nC})] Ensemble B limité aux épitaphes (\st{nEpit}) et aux textes mixtes, qui associent une commémoration et un acte d’achat (\st{nMixte}). Il exclut l’inscription non funéraire 321 et les \st{nActe} actes d’achat ou de réservation d’emplacement sans commémoration. C’est le dénominateur de toutes les analyses du formulaire, du vocabulaire religieux, des âges et des signes.
\end{description}

Les unités documentaires — monuments supposés distincts — ne forment pas un quatrième ensemble analytique : faute d’examen des supports, leur nombre n’est connu que par les identités que les bases établissent. Dans l’ensemble B, un seul cas réunit deux textes sur un même monument (313 et 314, ICUR IV 11117) ; les \st{nB} textes correspondent donc, au mieux de la documentation, à \st{nD} unités.

\begin{figure}[tbp]\centering\small
\setlength{\fboxsep}{5pt}
\fbox{\begin{minipage}{0.86\linewidth}\centering
Export EDCS « province Roma », 9 décembre 2020 (93\,058 fiches)\\[2pt]$\downarrow$ critères (a)–(c)\\[2pt]
\textbf{Ensemble A : \st{nA} notices} (identifiants stables)\\[2pt]
$\downarrow$ moins 3 doublons (106→120, 282→288, 329→336), 1 fragment joint (394→393), 1 texte hors période (98)\\[2pt]
\textbf{Ensemble B : \st{nB} textes distincts} (\st{nD} unités documentaires)\\[2pt]
$\downarrow$ moins 1 inscription non funéraire (321) et \st{nActe} actes d’achat\\[2pt]
\textbf{Ensemble C : \st{nC} textes funéraires} (\st{nEpit} épitaphes, \st{nMixte} textes mixtes)
\end{minipage}}
\caption{Constitution du corpus et ensembles analytiques.}\label{fig:corpus}
\end{figure}

\section{Doublons, fragments et rapprochements}

La recherche des doublons a combiné la similarité des textes, calculée entre notices d’années voisines, et le partage d’une référence d’édition. Aucune notice n’a été fusionnée sur la seule base de la similarité : chaque décision s’appuie sur une concordance publiée par une base savante (tableau~\ref{tab:doublons}).

\begin{table}[tbp]\centering\small
\caption{Identités documentaires et décisions.}\label{tab:doublons}
\begin{tabularx}{\linewidth}{@{}l l X@{}}\toprule
Notices & Décision & Preuve et remarques\\\midrule
329 / 336 & fusion dans 336 & L’EDB (EDB10126 = ICVR II 6050) cite l’AE 1987, 61, unique référence du \nnum{329} ; l’EDH (HD009698) concorde. Âge : \emph{L[X]V} (EDB), le X étant restitué.\\
106 / 120 & fusion dans 120 & L’EDH (HD030246) et l’EDB (EDB13167) rattachent l’AE 1905, 71 à ICUR II 6028. Leçon divergente du verbe : \emph{c(iu)esq(it)} / \emph{recessit}.\\
282 / 288 & fusion dans 288 & L’EDH (HD021839) et l’EDB (EDB11278) rattachent l’AE 1906, 86 à ICUR II 6043. Nom : \emph{An]iciae} (EDH) / \emph{Euticiae} (ICVR, EDB).\\
393 / 394 & un seul texte & L’EDB (EDB39590 = ICVR IV 11126) publie une fiche unique pour deux petits fragments (a, b) d’une même plaque ; elle ne lit ni \emph{Arc-} ni \emph{Hon-}. Un seul texte daté de 396.\\
313 / 314 & deux textes, un monument & L’EDB (EDB39579 = ICVR IV 11117) décrit une plaque portant deux dépositions successives : un défunt (389) et son frère (année restituée).\\
98 & hors période & \emph{Severin[us]} lu ; \emph{[Dagalaiphi]} restitué. L’EDB (CLRE p. 680) et l’EDH (HD028015) datent le texte de 461 ou 482.\\
270 / 274 & maintenus séparés & Textes réduits à une formule consulaire identique ; références ICUR distinctes.\\
\bottomrule\end{tabularx}
\end{table}

Ces identités ne sont pas établies sur les pierres. Elles reposent sur la concordance de deux bases indépendantes, qui citent les mêmes éditions ; c’est le niveau de preuve le plus élevé dont le volume dispose. Leur effet sur les mesures est faible : la proportion de textes funéraires portant \emph{in pace} passe de 41,9~\% dans l’ensemble initial de 468 textes funéraires à \st{ip_pct}~\% dans l’ensemble C (chapitre~\ref{ch:mesures}).

\section{Genre épigraphique et périmètre}

Les \st{nA} notices ne sont pas toutes des épitaphes. Quatre genres sont distingués, par une règle explicite — présence d’un verbe d’acquisition, présence ou absence d’une formule de commémoration — contrôlée à la lecture :
\begin{itemize}
\item les \emph{épitaphes}, qui commémorent un ou plusieurs défunts ;
\item les \emph{actes d’achat ou de réservation} (\nnums{\st{acteslist}}), qui attestent la propriété d’un emplacement sans commémorer de défunt ;
\item les \emph{textes mixtes}, qui combinent les deux fonctions ;
\item une \emph{inscription non funéraire} (\nnum{321}), qui commémore une colonne de la basilique de l’apôtre Paul sous l’évêque Sirice (390). Elle est conservée dans l’ensemble B et exclue de toutes les mesures funéraires.
\end{itemize}

Le texte 406 (\emph{Sussanna comparavit sibi memoriam quievit die VII Kal(endas) Augusta⸢s⸣}) est classé parmi les textes mixtes. Le verbe \emph{quievit} précède la date : celle-ci se rapporte au repos, c’est-à-dire au décès, de Susanna, et non à l’acquisition. Le classement inverse, comme acte d’achat, est testé au chapitre~\ref{ch:mesures} ; il ne modifie aucun résultat de manière significative.

\section{Trois notices à statut particulier}

\emph{Provenance douteuse (\nnum{439}).} L’épitaphe de Petulus, datée par le consulat de Stilichon (400), est identifiée comme CIL XIV 3896 = \emph{Inscriptiones Italiae} IV 1, 543, c’est-à-dire comme une inscription de Tibur. L’EDB la range parmi les inscriptions d’origine incertaine. Elle est conservée dans les ensembles, mais exclue de toute analyse topographique propre à Rome.

\emph{Datation prosopographique (\nnum{36}).} L’épitaphe de l’épouse de Clodius Adelfius, \emph{ex praefectis urbis}, ne porte aucun élément de date. L’année 351 est celle de la préfecture urbaine de son mari, c’est-à-dire un terme postérieur. La notice est exclue des séries chronologiques fines.

\emph{Date possiblement hors période (\nnum{52}).} Seules les lettres \emph{Da[} sont lues après \emph{conss(ulatu)}. Silvagni restitue Datianus (358) ; la CLRE, citée par l’EDB, rappelle que Dalmatius, consul en 333, commence aussi par DA. La notice est conservée sous l’année 358, avec une datation déclarée contestée.

\section{Ce que le corpus représente}

Le corpus n’est pas un échantillon de la population chrétienne de Rome, ni même de ses inscriptions funéraires. Il résulte d’une suite de sélections : la gravure d’une date consulaire, qui n’a concerné qu’une minorité des épitaphes ; la conservation et la découverte des pierres ; leur publication ; leur saisie dans une base ; enfin les critères d’extraction du présent travail. Chaque proportion donnée dans ce livre décrit cette documentation transmise. Elle ne devient une estimation de pratiques antiques qu’au prix d’hypothèses supplémentaires, qui sont discutées lorsqu’elles sont faites.
